
\chapter{Tensorgrad}

实现细节

\section{同构}
实际上存在“张量同构”（tensor isomorphism）的概念，但它基本上与图同构相同。

我们需要在代码的许多不同部分理解同构。

\subsection{在乘积中}
- 消去 / 合并乘积中的相同部分
这其实格外困难，因为必须收集构成同构子图的一组节点。
目前我们稍微取巧，只考虑乘积中的独立连通分量。

\begin{figure}[h]
\centering{
\def\svgwidth{\linewidth}
\import{figures/}{power.pdf_tex}
\caption{合并乘积中的相同部分}
\label{fig:power}
}
\end{figure}

这个问题基本上是：
\begin{enumerate}
   \item 给定一个多重图 $G$，其节点为 $V$、边为 $E$。
   \item 节点和边都有标签。
   \item 要找到 $V$ 的两个不相交子集 $V_1$ 和 $V_2$，使得由 $V_1$ 和 $V_2$ 诱导出的子图 $G_1$ 和 $G_2$ 同构。
      此外，在该同构下，$G_1$ 和 $G_2$ 中节点与边的标签相同。
\end{enumerate}

这个问题很可能是 NP-hard 的，但也许仍然存在比尝试所有子集的 $2^n$ 更快的算法。
特别是，我们可以修改 VF2 算法；该算法会迭代地尝试匹配 $G_1$ 和 $G_2$ 中的节点。
NetworkX 库已经有一个 GraphMatcher，可以搜索同构子图。
它也许可以扩展到我们的问题……
但老实说，我不知道我们是否真的想在最一般的情形下解决这个问题，因为它有点对应于对图做因式分解。
而我们并不做因式分解，就像我们不做逆分配律一样。

无论如何，显然我们需要能够按同构来比较节点和边。

此外，乘积中同构规范化的基本用例，只是根据各个部分计算乘积本身的规范形式。
这里我们方法的一部分，是把外部边转成节点，从而可以给它们着色。

\subsection{在求和中}

在判断 $A + B$ 是否等于 $2A$ 时，需要检查 $A$ 和 $B$ 是否同构。
但我们还必须在当前边重命名下完成这件事。
这就是为什么不能直接把 $A + A^T$ 变成 $2A$。

我的代码中实际的工作方式是
\begin{lstlisting}
def key_fn(t: Tensor):
    # Align tensor edges to have the same order, using Sum's order as reference.
    canons = [t.canonical_edge_names[t.edges.index(e)] for e in self.edges]
    return hash((t.canon,) + tuple(canons))

ws_tensors = TensorDict(key_fn=key_fn, default_fn=int)
for w, t in zip(weights, tensors):
    ws_tensors[t] += w
ws_tensors = [(w, t) for t, w in ws_tensors.items()]
\end{lstlisting}
这表示我用于 hash 的内容，是张量的规范形式，加上按照求和中边顺序排列的边的规范形式。
这些基本上就是轨道；这意味着如果被求和的张量具有某种对称性，我们就允许“翻转”它，从而使各项同构。

在 “compute canonical” 方法中，我们做的事情大致相同，但还会包含权重。
\begin{lstlisting}
def _compute_canonical(self):
    hashes = []
    for e in self.edges:
        canons = [t.canonical_edge_names[t.edges.index(e)] for t in self.tensors]
        hashes.append(hash(("Sum",) + tuple(sorted(zip(self.weights, canons)))))
    base = hash(("Sum", len(self.tensors)))
    hashes = [hash((base, h)) for h in hashes]
    return base, hashes
\end{lstlisting}

将来我们想改用对称群。
求和的对称群会是什么？
它是各个求和项对称群乘积的对角部分。
怎样才能找到这个群的生成元？
也许应该直接构造某个联合图，然后寻找该图的自同构。

另一种选择是使用 sympy。
目前尚不清楚这个问题是否可以在多项式时间内求解。我认为 Babai 证明了它是准多项式的，但该算法并不实用。顺带一提，子群交、元素中心化子以及 $\{1, \dots, n\}$ 子集稳定子这些问题已经被证明（由 Eugene Luks 证明）在多项式意义下等价。


实际上，构造一个图并使用 nauty 是个很好的想法，因为它能够检测出 $A+A^T$ 是对称的。
仅取各个求和项自同构群的交集并不能发现这一点。

另一种选择是把求和转换成一个函数……
但不，这有点奇怪。
那会要求我支持任意数量输入的函数，而目前还不是这样。



\subsection{在求值中}
在对张量求值时，可以查看该张量的图，并判断它是否同构于先前已经求值过的某个张量。
这是一个并不真正需要规范形式的例子；近似 hash 加上 vf2 就足够了。
还要注意，在这种情况下我们并不关心边的重命名，因为可以在返回张量之前直接重命名边。
例如，如果已经求值过 $A$，就可以很容易用它得到 $A^T$。

\subsection{在变量中}
在变量中，我们把变量名包含进 hash。
基本上，我们假设同名变量指向同一份数据。
\begin{lstlisting}
   base = hash(("Variable", self.name))
   return base, [hash((base, e)) for e in self.original_edges]
\end{lstlisting}
对于原始规范边名，我们使用重命名前的边名。
这意味着在 $A^T$ 的情形下，它会与 $A$ 具有相同的 hash。
但因为它已经被重命名，Sum 中的 \texttt{t.index} 调用会翻转这些边。

可以设想让变量接收一个自同构群作为参数，这样就能定义具有不同对称性的变量。
例如一个对称矩阵 $A$，此时 $A+A^T$ 实际上就是 $2A$。

\subsection{在常量中}
在计算 \texttt{Zero} 或 \texttt{Copy} 这类常量的规范形式时，我们并不关心边名。
我猜这是因为我们使用的常量都是最大对称的？
目前我们会包含常量的 \texttt{tag}；如果它来自某个变量，那么这个 tag 就是该变量的 hash。

\subsection{在函数中}
一个问题是，原始名称通常是函数定义的一部分，
而求导新增的边往往是根据上下文自动生成的，
因此它们其实不应该成为规范形式的一部分。

与 Sum 不同，这里我们不会对规范名排序，因为输入顺序是重要的。

也许应该允许函数变换对称群？
例如，如果有一个函数接收对称矩阵并返回对称矩阵，那么我们应当能使用输入的对称群来化简输出。

\subsection{在导数中}
我们所做的只是对张量和 wrt 做 hash。
然后为导数添加新边。

\subsection{其他}
对于某些张量，可能存在并非等价关系的边维度关系。
例如，一个 flatten 张量会让 “out” 边维度等于各个 “in” 边维度的乘积。

在之前的版本中，我让每个张量都注册一个“回调”（callback）函数。
每当某个边维度“变得可用”时，该张量就有机会发出新的边维度。
然而，这会给每个张量的实现增加很多工作，而且现有张量都不需要它。

\section{重命名}
这是代码中一个重要部分。

\section{求值}

求值中的一个重要部分是确定每条边的维度。
为此，我基本上会构造张量的完整图，并使用一个名为
\texttt{edge\_equivalences} 的函数，它给出一列元组 $((t_1, e_1), (t_2, e_2))$，
表示张量 $t_1$ 的边 $e_1$ 等价于张量 $t_2$ 的边 $e_2$。
注意，同一个边名可能在图中出现多次，因此也需要跟踪张量本身。

对于变量，由于用户以变量形式给出边维度，因此跟踪重命名后的边名很重要：
\begin{lstlisting}
for e1, e2 in zip(self.original_edges, self.edges):
     yield (self, e1), (self, e2)
\end{lstlisting}

对于常量，可能存在一些基于该常量来源张量的等价关系。
\begin{lstlisting}
def edge_equivalences(self):
    if self.link is not None:
        yield from self.link.edge_equivalences()
        for e in self.link.edges:
            if e in self.edges:
                yield (self, e), (self.link, e)
\end{lstlisting}

对于复制张量，所有边都是等价的：
\begin{lstlisting}
def edge_equivalences(self):
    yield from super().edge_equivalences()
    for e in self.edges[1:]:
        yield (self, self.edges[0]), (self, e)
\end{lstlisting}

对于函数，我们实际上无法对函数自身的边 \texttt{(self.edges\_out)} 说什么，
但至少可以对广播边说一些事情。
\begin{lstlisting}
for t, *inner_edges in self.inputs:
    yield from t.edge_equivalences()
    for e in t.edges:
        if e not in inner_edges:
            yield (t, e), (self, e)
\end{lstlisting}
也许还可以说同名输入边是等价的？

对于乘积，我们查看每条边 $(t_1, e, t_2)$，并产生 $(t_1, e), (t_2, e)$。
不过对于自由边 $(t, e)$，我们会把它们与自身匹配，即 $(t, e), (self, e)$。
\begin{lstlisting}
def edge_equivalences(self):
    pairs = defaultdict(list)
    for t in self.tensors:
        yield from t.edge_equivalences()
        for e in t.edges:
            pairs[e].append(t)
    for e, ts in pairs.items():
        if len(ts) == 1:
            yield (self, e), (ts[0], e)
        else:
            t1, t2 = ts
            yield (t1, e), (t2, e)
\end{lstlisting}

类似地，对于求和，一切都只是与自身匹配：
\begin{lstlisting}
def edge_equivalences(self):
    for t in self.tensors:
        yield from t.edge_equivalences()
        for e in t.edges:
            yield (t, e), (self, e)
\end{lstlisting}

最后，我们使用 BFS 将边维度从变量（由用户给定）传播到图的其余部分。

为什么非变量也需要知道边维度？
主要是因为复制张量：我们会用它们表示超边，并且必须构造它们。
如果我们不用 copy 而是更高效地计算超边，能否摆脱这一点？
有时也会出现“分离的”（detached）copy……

此外，另一种想法是真正构造完整图。
我原本认为这不可行，因为 Sum 并不真的是图。
但也许借助使用 nauty 的新方法，我们实际上可以做到这一点。

\subsection{乘积}
我们只需对乘积中的张量求值，并把它们交给 einsum。


\section{化简规则}
这类规则有很多。

大多数情况下，我们可以在一次深度优先遍历中完成所有事情，
但有少数时候需要多轮遍历。
这可以通过 full-simplify 方法完成；它会反复调用 simplify，直到不再发生变化。
